%%%%%%%%%%%%%%%%%%%%%%%%%%%%%%%%%%%%%%%%%%%%%%%%%%%%%%%%%%%%%%%%%%%%%%%%%%%%%%%%%%%%%%%%%%%%%%%%%%%%
%%% +++---~~~~ Begin: main document ~~~---+++ %%%%%%%%%%%%%%%%%%%%%%%%%%%%%%%%%%%%%%%%%%%%%%%%%%%%%%
%%%%%%%%%%%%%%%%%%%%%%%%%%%%%%%%%%%%%%%%%%%%%%%%%%%%%%%%%%%%%%%%%%%%%%%%%%%%%%%%%%%%%%%%%%%%%%%%%%%%
%%% Author:         Fabian Roos, fabian@froos.de
%%% Datum:          01.02.2016
%%% Document:       Example Thesis Document
%%% Information:    Compile document with './compile_main_example.sh' or alternatively
%%%                 -> 'pdflatex -shell-escape main_example.tex'
%%%                 -> 'makeglossaries main_example'
%%%                 -> 'bibtex main_example'
%%%                 -> 'pdflatex -shell-escape main_example.tex'
%%%                 -> 'pdflatex -shell-escape main_example.tex'
%%%%%%%%%%%%%%%%%%%%%%%%%%%%%%%%%%%%%%%%%%%%%%%%%%%%%%%%%%%%%%%%%%%%%%%%%%%%%%%%%%%%%%%%%%%%%%%%%%%%



%%% +++---~~~ Load Header: Settings for the Document ~~~---+++ %%%%%%%%%%%%%%%%%%%%%%%%%%%%%%%%%%%%%
\newcommand{\LANGUAGE}{1}                       % 1 = german, 0 = english
\newcommand{\SHOWDATETIME}{1}                   % 1 = show or 0 = hide the date and time in footer
\newcommand{\TYPE}{1}                           % 1 = master, 0 = bachelor
%%%%%%%%%%%%%%%%%%%%%%%%%%%%%%%%%%%%%%%%%%%%%%%%%%%%%%%%%%%%%%%%%%%%%%%%%%%%%%%%%%%%%%%%%%%%%%%%%%%%
%%% +++---~~~~ start: settings document template: header.tex~~~---+++ %%%%%%%%%%%%%%%%%%%%%%%%%%%%%%
%%%%%%%%%%%%%%%%%%%%%%%%%%%%%%%%%%%%%%%%%%%%%%%%%%%%%%%%%%%%%%%%%%%%%%%%%%%%%%%%%%%%%%%%%%%%%%%%%%%%
%%% Autor:      Fabian Roos, fabian@froos.de
%%% Datum:      01.02.2016 (Nicht Einstellungsänderung, sondern Inhalt!)
%%% Dokument:   Vorlage Masterarbeit (/ Praktikumsbericht) (Vim: kein 'gg=G'!)
%%% Version:    0.1: Erster Entwurf, verwendet für Praktikumsbericht
%%%             0.2: Umstellung auf Paket 'tikzscale'
%%%	            0.3: Restrukturierung der Grafik-Pakete und -Optionen, final für Masterarbeit
%%%             0.4: Einbindung der MWT-Optionen, geometry-Qualmöglichkeit, LaTeX_RS-Version
%%%             0.5: Englische Version für babel, siunitx-Optionen
%%%             0.6: Befehl für Referenzen auf Legenden in PGFPlots
%%%             0.7: Unterstützung für Inkscape-Dateien
%%%             0.8: Flag für Sprache, myunits
%%%             0.9: prepared biber bibliography support
%%%%%%%%%%%%%%%%%%%%%%%%%%%%%%%%%%%%%%%%%%%%%%%%%%%%%%%%%%%%%%%%%%%%%%%%%%%%%%%%%%%%%%%%%%%%%%%%%%%%


%%% +++---~~~ begin: documentclass definitions ~~~---+++ %%%%%%%%%%%%%%%%%%%%%%%%%%%%%%%%%%%%%%%%%%%
% the documentation: http://mirrors.ctan.org/macros/latex/contrib/koma-script/doc/scrguide.pdf
\documentclass[%
% ~~~~~~~~~ Document ~~~~~~~~~~~~~~~~~~~~~~~~~~~~~~~~~~~~~~~~~~~~~~~~~~~~~~~~~~~~~~~~~~~~~~~~~~~~~~~
    % final,    % final document, default option
    % draft,    % show 'overfull hboxes' to see to large elements, no graphics are shown!
% ~~~~~~~~~ font settings ~~~~~~~~~~~~~~~~~~~~~~~~~~~~~~~~~~~~~~~~~~~~~~~~~~~~~~~~~~~~~~~~~~~~~~~~~~
    % fontsize=11pt,    % default value for scrbook, scrreprt, and scrartcl is 11pt, for PHD uses 9pt?
    % headings=normal,  % font size for the headings, which are set large in the scrclass
                        % default should be =big, but =small or =normal are possible as well
% ~~~~~~~~~ Layout ~~~~~~~~~~~~~~~~~~~~~~~~~~~~~~~~~~~~~~~~~~~~~~~~~~~~~~~~~~~~~~~~~~~~~~~~~~~~~~~~~
    paper=a4,	% Einstellung Papiergröße, Standard bei KOMA-Script, s. Handbuch S. 47
    pagesize,	% Schreibt die Dateigröße ins PDF, sollte immer gesetzt sein, s. Handb. S.47
    % landscape,% Landscape-Format
    % oneside,	% Formatierung für einseitiger Druck, Standardoption
    twoside,	% Formatierung für doppelseitiger Druck
    % BCOR=0.78mm,  % Wie viel Platz geht pro Seite für die Bindung verloren? S. 32
    % onecolumn	% Einspaltig, Standard
    % twocolumn,% Zweispaltig
    open=right,	% Wo soll neues Kapitel beginnen? Nur bei scrreprt sinnvoll! S. 89
                % 'any':	Auf jeder neuen Seite
                % 'right':	Auf jeder rechten Seite, fügt ggf. leere Seiten ein!
                %           Ende einer Seite immer auf einer Ebene (\flushbottom wird gesetzt!), S. 29
                %           Dies kann dazu führen, dass es mitunter große Absatzabstände gibt. Daher
                %           ggf. ans Ende einer Seite ein \clearpage http://www.komascript.de/flushbottom
                % 'left':	Auf jeder linken Seite, fügt ggf. leere Seiten ein!
    % parskip=half,     % Absatzabstand /- einzug.
                        % Standard: false, kein Abstand, sondern 1em-Einzug. S.77
    % DIV=calc,   % only for a4 papers default values for the DIV value are defined, if other paper
                % format is chosen, typearea claculates a good value.
                % setting the calculation overwrites the default, use with caution
% ~~~~~~~~ Titelseite ~~~~~~~~~~~~~~~~~~~~~~~~~~~~~~~~~~~~~~~~~~~~~~~~~~~~~~~~~~~~~~~~~~~~~~~~~~~~~~
    % titelpage=false,  % Titelkopf wird gesetzt, Standard bei scrarctl
    % titelpage,        % Titelei in Form von Titelseite, Standard bei scrbook, scrreprt
% ~~~~~~~~ Kopf-/Fuß-Zeile ~~~~~~~~~~~~~~~~~~~~~~~~~~~~~~~~~~~~~~~~~~~~~~~~~~~~~~~~~~~~~~~~~~~~~~~~~
    % headinclude,      % Kopfzeile mit bei der Satzspiegelberechnung beachtet werden? Handb. S. 41
    headsepline,        % Unter Kolumnentitel eine separate Linie setzen, siehe S. 70
    %footsepline,       % Über dem Fuß eine separate Linie setzen, siehe S. 70
% ~~~~~~~~ Verzeichnisse ~~~~~~~~~~~~~~~~~~~~~~~~~~~~~~~~~~~~~~~~~~~~~~~~~~~~~~~~~~~~~~~~~~~~~~~~~~~
    % abstract=true,    % Bei der \begin{abstract}-Umgebung schreibt dies "Zusammenfassung" darüber
    toc=bibliography,   % Literaturverzeichnis ins Inhaltsverzeichnis, S. 73
    % toc=bibliographynumbered,	% Dieses Mal numeriert, S. 73
    % toc=indent,       % Inhaltsverzeichnis hierarchisch, alternativ =graduated, Standard, S. 73
    % toc=flat,         % auch =left, Inhatlsverzeichnis tabellarisch: Gliederungsnummern erste
                        % Spalte, Überschriften zweite, Seitenzahlen dritte Spalte, siehe S. 73
    % toc=index,        % Stichwortverzeichnis in den Index, auch =idx, siehe S. 73
    toc=listof,         % Gleitumgebungen (Abbildungs-, Tabellenverzeichnis) ins Inhaltsverz. S. 72
    % ?                 % weitere Optionen siehe S. 72 und S. 135
    numbers=noenddot,   % Überschritfsnummern ohne Punkt am Ende, alt pointlessnumbers, S. 98
% ~~~~~~~~ Mathematik ~~~~~~~~~~~~~~~~~~~~~~~~~~~~~~~~~~~~~~~~~~~~~~~~~~~~~~~~~~~~~~~~~~~~~~~~~~~~~~
    intlimits,          % place the limits above and below instead of right
    % leqno,            % Formelnummer auf der linken Seite anzeigen, S. 120, amsmath: S. 3
    % fleqn,            % Formeln linksbündig ausrichten, S. 120, amsmath: S. 3
]{scrreprt}             % Komascript-Klassen: scrartcl, scrreprt, scrbook
%%% +++---~~~ Ende: documentclass-Definitionen ~~~---+++  %%%%%%%%%%%%%%%%%%%%%%%%%%%%%%%%%%%%%%%%%%


%%% +++---~~~ Kopfzeile aktivieren ~~~---+++ %%%%%%%%%%%%%%%%%%%%%%%%%%%%%%%%%%%%%%%%%%%%%%%%%%%%%%%
\pagestyle{headings}	% 'plain' ist leere Kopf- und Fußzeile, nur Seitenzahl, Standard, S. 76
                        % 'headings' gibt zusätzlich lebenden Kolumnentitel aus
                        % Besondere Kopfzeilen für u.a. Titelseite, Indexseite: S. 79

% MWT: Einkommentieren: Nur für Datum in Fußzeile? Dann für final auskommentieren?
\usepackage{scrpage2}   % Nur für folgende Definitionen verwenden
\usepackage{scrtime}    % Verschiedene Zeitformate
\pagestyle{scrheadings}
\clearscrheadfoot%      % delete all entries
                        % define header and footer fields
\automark[section]{chapter} % \automark[rechte Seite]{linke Seite} S. 225
\ohead{\headmark}
\ofoot[\pagemark]{\pagemark}
\ifnum\SHOWDATETIME=1
    \ifoot[\tiny\today, \thistime]{\tiny\today, \thistime}
\fi


%%% +++---~~~ Verzeichnis beeinflussen ~~~̣---+++ %%%%%%%%%%%%%%%%%%%%%%%%%%%%%%%%%%%%%%%%%%%%%%%%%%
                        % -1 \part | 0 \chapter | 1 \section | 2 \subsection (Standard)
                        % 3 \subsubsection | 4 paragraph | 5 subparagraph
% \setcounter{tocdepth}{3}      % Bis zu welcher Gliederungsebene im Inhaltsverzeichnis? Siehe S. 70
% \setcounter{secnumdepth}{3}   % Bis zu welcher Gliederungsebene Nummer vor Titel anzeigen? S. 106
                        % Wird ein Kapitel per '\chapter*{Name}' ohne Kapitelnummer erstellt, so
                        % wird es auch nicht im Inhaltsverzeichnis angezeigt.
                        % -> Setze vor die Kapiteldefinition folgende Zeile und das Kapitel taucht
                        %    ToC auf
                        % \addcontentsline{toc}{chapter}{Symbole}
                        % https://tex.stackexchange.com/questions/14379

% MWT: Einkommentieren
% Definition sprachabhängiger Bezeichner, S. 297
\ifcase \LANGUAGE
    \renewcaptionname{english}{\figurename}{Fig.}
    \renewcaptionname{english}{\tablename}{Tab.}
\else
    \renewcaptionname{ngerman}{\figurename}{Bild}
    \renewcaptionname{ngerman}{\tablename}{Tab.}
\fi

% \numberwithin{equation}{section} % (Chapter.Section.Number) statt (Chapter.Number)


%%% +++---~~~ Seitennummerierung beeinflussen ~~~̣---+++ %%%%%%%%%%%%%%%%%%%%%%%%%%%%%%%%%%%%%%%%%%%%
% \pagenumbering{?}     % Setzt die Seitennummerierung auf '1' zurück, wählt Zählweise. Siehe S. 80
                        % arabic (8) / roman (viii) / Roman (VIII) / alph (h) / Alph (H)
                        % Im Dokument an die Stelle setzen, an der die neue Numerierung beginnen
                        % soll
                        % Ggf. muss die Seite erst beendet werden mit '\clearpage' bzw.
                        % '\cleardoublepage'


%%% +++---~~~ Anfang: Pakete laden ~~~---+++ %%%%%%%%%%%%%%%%%%%%%%%%%%%%%%%%%%%%%%%%%%%%%%%%%%%%%%%
% ~~~~~~~~ Standard-Pakete ~~~~~~~~~~~~~~~~~~~~~~~~~~~~~~~~~~~~~~~~~~~~~~~~~~~~~~~~~~~~~~~~~~~~~~~~~
% Für die drei folgenden Pakete siehe auch: http://de.wikipedia.org/wiki/LaTeX#Zeichenkodierung

% <!-- Pakete für compilieren mit pdflatex
\usepackage[utf8]{inputenc}	% Zeichencodierung
\usepackage[T1]{fontenc}	% Schriftart, für Umlaute
% -->

% <!-- Pakete für compilieren mit lualatex
% \usepackage{luainputenc}
% -->

\usepackage{lmodern}        % Hochwertige Fonts: http://tex.stackexchange.com/questions/147194/
\ifcase \LANGUAGE           % load more than one language, active one last
                            % http://tex.stackexchange.com/questions/20987/
    \usepackage[ngerman,english]{babel} % english
\else
    \usepackage[english,ngerman]{babel}	% Für deutsche Begriffe (Inhaltsver., Kapitel) und Worttrennung
                                        % Alte Rechtschreibung: german
                                        % Neue Rechtschreibung: ngerman
                                        % Quelle: http://tex.stackexchange.com/questions/67549/
\fi
% \usepackage[]{csquotes}

% ~~~~~~~~ Grafik-Pakete ~~~~~~~~~~~~~~~~~~~~~~~~~~~~~~~~~~~~~~~~~~~~~~~~~~~~~~~~~~~~~~~~~~~~~~~~~~~
% \usepackage{pdfpages}     % Für das Einbinden von PDF-Dateien als ganze Seite (nicht als Bild)
                            % \includepdf[width=\linewidth]{x.pdf}
\usepackage{graphicx}       % Grafiken: http://en.wikibooks.org/wiki/LaTeX/Importing_Graphics
\usepackage[dvipsnames]{xcolor} % Lade ein definiertes Set an Farben
\usepackage[%               % Beides MWT-Optionen
    hang,
    labelfont = bf,
]{caption}                  % Wird für 'subcaption' benötigt in alten TexLive-Versionen
\usepackage{subcaption}     % Statt veraltete Pakete 'subfigure' / 'subfig'. Siehe Sündenregister
% \usepackage{wrapfig}      % Falls Grafiken von Text umflossen werden sollen, eine Möglichkeit!
                            % Für Grafiken muss Paket graphicx geladen werden
% \usepackage{morefloats}   % Mehr als 18 unbearbeitete floats: http://ctan.org/pkg/morefloats
% \usepackage{flafter}      % MWT: Floats nach Referenz setzen
% \usepackage{placeins}     % MWT: Alle unbearbeiteten Floats anzeigen.
                            % Erlaubt mit \Floatbarrier das Ausgeben aller bisherigen Floats
                            % Achtung: Kann zu teilweise leeren Seiten führen, nur mit Bedacht
                            % verwenden. Besser wäre es mit [!t] oder [!b] zu arbeiten.
\usepackage{adjustbox}      % Paket für die Skalierung von Flussdiagrammen
% \newsavebox{\measurebox}  % verwendet für die Platzierung dreier Bilder nebeneinander

% <!-- Folgende Pakete auch in export.tex
% TikZ
\usepackage{tikz}           % Für Tikz-Grafiken
\usetikzlibrary{%
    external,               % Bibliothek für das externalisieren von TikZ
    calc,                   % Für die Berechnung von Parametern
    shapes.geometric,       % Flussdiagramme
    arrows,                 % Flussdiagramme
    fit,                    % Flussdiagramme: Boxen um Nodes zeichnen
    trees,                  % für 'circular growth' bei den Baumdiagrammen
    % snakes,               % superseded, 'decorations.pathmorphing' verwenden
    decorations.pathmorphing,
}
\usepackage{tkz-euclide}    % Für die Einfache Kennzeichnung von Winkeln
\usetkzobj{all}             % Option für 'tkz-euclide'

% PGF
\usepgflibrary{%
    plotmarks,              % Scatterplots: Kreise statt Linien
}

% PGFPlots (auch in 'export*' benötigt - standalone-Paket-Datei)
\usepackage{pgfplots}       % Paket zum Plotten laden
\usepgfplotslibrary{%
    external,               % Bibliothek für das externalisieren von PGFPlots
    statistics,             % Box Plots
    groupplots,             % Groupplots
    smithchart,             % Smith Charts
    polar,                  % polar plots
    fillbetween,            % fill area with colour
}
\pgfplotsset{compat=newest} % Verwende immer neueste Version und kein Kompatibilitätsmodus
% \pgfplotsset{plot coordinates/math parser=false} % Für matlab2TikZ

% Skalierung von TikZ / PGF
\usepackage{tikzscale}      % Einbinden von TikZ/ PGFPlots zum skalieren per
                            % \includegraphics[width=\linewidth,height=0.9\textheight]{gfx.tikz}
                            % Höhe bei TikZ sinnlos, bei PGFPlots für Achsenverhältnis!
                            % Grafik muss *.tikz sein, Fehler bei *.tex

\usepackage[%
    european,
  ]{circuitikz}             % Für Schaltungen

% Define the Uni Ulm colors
\definecolor{uniblue}{RGB}{125,154,170}
\definecolor{unired}{RGB}{163,38,56}
%
\tikzset{%
    % arrow shape: stealth, latex
    >=stealth,
    %
    %% flow charts
    % -> rectangle blue
    fcrecb/.style={rectangle, rounded corners, align=center, draw=blue, fill=blue!15},
    % -> ellipse red
    fcelr/.style={ellipse, align = center, draw=red, fill=red!20},
    % -> use in document as
    % \tikzset{%
    %   fcblock/.style={fcrecb, minimum width=4cm, minimum height=1.25cm},
    %   fcdecision/.style={fcelr, minimum width=6cm, minimum height=1.25cm},
    % }
}
%
\pgfplotscreateplotcyclelist{roos}{%
  % Usage: Add the following in th axis environment
  % cycle list name = roos,
  % no markers,
  % mark repeat={4},  % Markers only at every nth Datapoint
  % \addplot+
  {blue, solid, mark=x, mark options={solid}},
  {red, densely dashed, mark=o, mark options={solid}},
  {orange, densely dashdotted, mark=triangle, mark options={solid}},
  {black, densely dotted, mark=square, mark options={solid}},
  {brown, loosely dashed, mark=Mercedes star, mark options={solid}},
}
%
\pgfplotscreateplotcyclelist{roosthick}{%
  % Usage: Add the following in th axis environment
  % cycle list name = roos,
  % no markers,
  % mark repeat={4},  % Markers only at every nth Datapoint
  % \addplot+
  {blue, solid, thick, mark=x, mark options={solid}},
  {red, densely dashed, thick, mark=o, mark options={solid}},
  {orange, densely dashdotted, thick, mark=triangle, mark options={solid}},
  {black, densely dotted, thick, mark=square, mark options={solid}},
  {brown, loosely dashed, thick, mark=Mercedes star, mark options={solid}},
}
%
\ifcase \LANGUAGE
\else
% For german documents the decimal seperator is a comma and not a dot. Set this globally
% http://tex.stackexchange.com/a/247423
\pgfplotsset{every linear axis/.append style={%
    /pgf/number format/use comma,
    /pgf/number format/1000 sep={},
  }
}
\fi
%
% Define the parula colormap
\pgfplotsset{% http://tex.stackexchange.com/questions/228197/
    colormap={parula}{%
        rgb = (0.2081,0.1663,0.5292)
        rgb = (0.2116,0.1898,0.5777)
        rgb = (0.2123,0.2138,0.627)
        rgb = (0.2081,0.2386,0.6771)
        rgb = (0.1959,0.2645,0.7279)
        rgb = (0.1707,0.2919,0.7792)
        rgb = (0.1253,0.3242,0.8303)
        rgb = (0.0591,0.3598,0.8683)
        rgb = (0.0117,0.3875,0.882)
        rgb = (0.006,0.4086,0.8828)
        rgb = (0.0165,0.4266,0.8786)
        rgb = (0.0329,0.443,0.872)
        rgb = (0.0498,0.4586,0.8641)
        rgb = (0.0629,0.4737,0.8554)
        rgb = (0.0723,0.4887,0.8467)
        rgb = (0.0779,0.504,0.8384)
        rgb = (0.0793,0.52,0.8312)
        rgb = (0.0749,0.5375,0.8263)
        rgb = (0.0641,0.557,0.824)
        rgb = (0.0488,0.5772,0.8228)
        rgb = (0.0343,0.5966,0.8199)
        rgb = (0.0265,0.6137,0.8135)
        rgb = (0.0239,0.6287,0.8038)
        rgb = (0.0231,0.6418,0.7913)
        rgb = (0.0228,0.6535,0.7768)
        rgb = (0.0267,0.6642,0.7607)
        rgb = (0.0384,0.6743,0.7436)
        rgb = (0.059,0.6838,0.7254)
        rgb = (0.0843,0.6928,0.7062)
        rgb = (0.1133,0.7015,0.6859)
        rgb = (0.1453,0.7098,0.6646)
        rgb = (0.1801,0.7177,0.6424)
        rgb = (0.2178,0.725,0.6193)
        rgb = (0.2586,0.7317,0.5954)
        rgb = (0.3022,0.7376,0.5712)
        rgb = (0.3482,0.7424,0.5473)
        rgb = (0.3953,0.7459,0.5244)
        rgb = (0.442,0.7481,0.5033)
        rgb = (0.4871,0.7491,0.484)
        rgb = (0.53,0.7491,0.4661)
        rgb = (0.5709,0.7485,0.4494)
        rgb = (0.6099,0.7473,0.4337)
        rgb = (0.6473,0.7456,0.4188)
        rgb = (0.6834,0.7435,0.4044)
        rgb = (0.7184,0.7411,0.3905)
        rgb = (0.7525,0.7384,0.3768)
        rgb = (0.7858,0.7356,0.3633)
        rgb = (0.8185,0.7327,0.3498)
        rgb = (0.8507,0.7299,0.336)
        rgb = (0.8824,0.7274,0.3217)
        rgb = (0.9139,0.7258,0.3063)
        rgb = (0.945,0.7261,0.2886)
        rgb = (0.9739,0.7314,0.2666)
        rgb = (0.9938,0.7455,0.2403)
        rgb = (0.999,0.7653,0.2164)
        rgb = (0.9955,0.7861,0.1967)
        rgb = (0.988,0.8066,0.1794)
        rgb = (0.9789,0.8271,0.1633)
        rgb = (0.9697,0.8481,0.1475)
        rgb = (0.9626,0.8705,0.1309)
        rgb = (0.9589,0.8949,0.1132)
        rgb = (0.9598,0.9218,0.0948)
        rgb = (0.9661,0.9514,0.0755)
        rgb = (0.9763,0.9831,0.0538)
    }
}
         % Definitionen für Flussdiagrame (Prefix 'fc'), Plotstyles, uvm.
% Ende der Pakete für export.tex -->

% after a \addplot in a PGFPlots figure a label \label{leg:snr} can be defined and referenced with
% \protect\plotref{leg:snr}
% source: http://tex.stackexchange.com/questions/65494/
% problem: warnings of Multiple defined labels: tex.stackexchange.com/questions/277583/
\newcommand{\plotref}[1]{\tikzexternaldisable\ref{#1}\tikzexternalenable}

% Sollte 'draft' gesetzt werden, so zeige die Boxen in der korrekten Größe an
% Quelle: http://tex.stackexchange.com/questions/47444/
\pgfkeys{/pgf/images/include external/.code=\includegraphics{#1}}

% Optimierung: Aktiviere 'externalize' (Setzen mit --shell-escape)
% \tikzexternalize[%
%     optimize command away=\includepdf, % Eingebundene PDFs nicht externalisieren
%                             % https://tex.stackexchange.com/questions/135504/
%     prefix=tikz/,           % Schreibe externalisierte PDF-Dateien in verzeichnis 'tikz'
% ]

% use lualatex for the externalization process, change inputenc or errors occur
% -> comment in only if you know what you do!
% \tikzset{external/system call={lualatex \tikzexternalcheckshellescape -halt-on-error -interaction=batchmode -jobname "\image" "\texsource"}}

% <!-- Include Inkscape files: Start
% \RequirePackage{scrlfile}                           % prevent the subfig package fromn being loaded
% \PreventPackageFromLoading{subfig}                  % because of known incompatibility with subcaption
%                                                     % use this until the svg package is updated
%                                                     % https://tex.stackexchange.com/questions/291929/
% % \expandafter\def\csname ver@subfig.sty\endcsname{}  % old, not working with Tex Live 2016
%                                                     % http://tex.stackexchange.com/questions/213273/
% \usepackage{svg}                                    % use this package to directly inser svg (!)
%                                                     % files into the document. There is no need for
%                                                     % the pdf_tex approach!
%                                                     % have a look at the README for further
%                                                     % information
% \setsvg{inkscape = inkscape -z -D}                  % with this only the drawing instead of the full
%                                                     % page is included
% \setsvg{svgpath = gfx/}                             % specify the folder where the svg files are
%                                                     % placed
% Include Inkscape files: End -->

\usepackage{xifthen}

% ~~~~~~~~ Mathemathik-Pakete ~~~~~~~~~~~~~~~~~~~~~~~~~~~~~~~~~~~~~~~~~~~~~~~~~~~~~~~~~~~~~~~~~~~~~~
\usepackage{amsmath}        % Für Mathematikumgebungen wie equation, align (& vor =, S. 3)
                            % Formeln referenzieren mit \eqref, siehe S. 10
                            % beinhaltet laut Dokumentation S. 2:
                            % -> amstext: \text innerhalb von Gleichungen
                            % -> amsopn: Neue Formeln mit \DeclareMathOperator
                            % -> amsbsy: Rückwärtskompatibilität
                            % Bei Problemen \usepackage{mathtools}
\usepackage{amssymb}		% Für besondere Zeichen, etwa \dashrightarrow, \mathbb{R}, lädt
                            % -> amsfont: Extra Fonts und Symbole
                            % http://de.wikibooks.org/wiki/LaTeX-Kompendium:_amssymb
\usepackage{trfsigns}       % Für \fourier also o-o
\usepackage{esint}          % special integrals, e.g. \oiin
\ifcase \LANGUAGE
    \usepackage[
        exponent-product={\cdot},       % Produktzeichen bei 1 \cdot 10^-5 ist Standard x (Version 2)
        per-mode=fraction,              % Ohne diese Angabe folgt m s^-1 statt als Bruch
        detect-all,
    ]{siunitx}
\else
    \usepackage[%                       % SIunits ist veraltet, siehe LaTeX2e-Sündenregister S. 15
                                        % Wie soll das Trennzeichen bei Zahlen aussehen?
        % decimalsymbol=comma,	        % Dezimaltrenner (Version 1)
        output-decimal-marker={,},      % Dezimaltrenner (Version 2)
        % expproduct = {\cdot},         % Produktzeichen bei 1 \cdot 10^-5 ist Standard x (Version 1)
        exponent-product={\cdot},       % Produktzeichen bei 1 \cdot 10^-5 ist Standard x (Version 2)
        per-mode=fraction,              % Ohne diese Angabe folgt m s^-1 statt als Bruch
        range-phrase={--},              % Für \SIrange{1}{10}{einheit} statt to ein -
        % range-phrase={ bis },         % Für \SIrange{1}{10}{einheit} statt to ein bis
        %range-units=single,            % Damit nicht nach jedem eine Einheit erscheint
        list-final-separator={ und },   % Für \SIlist{1;2;3}{\einheit}: Letzte Trennung wie?
        list-pair-separator={ und },    % Für \SIlist{1;2}{\einheit}: Trennung in der Mitte wie?
        % list-units=brackets,          % Damit bei SIlist (1;2;3) Einheit erscheint statt 1Einheit, 2Einheit
        detect-all,                     % Damit in strong-Umgebungen auch so dargstellt. Ggf. andere Option?
    ]{siunitx}                          % Beispiel: \SI{1,1}{\metre} kein '{,}' nötig
\fi

\setlength\delimitershortfall{-0.5pt}   % Damit Abstand zwischen mehreren Klammern bei Verwendung
                                        % von \left und \right eingefügt wird. Frank Bögelsack

\allowdisplaybreaks         % Damit in Gleichungen wie align Umbrüche erlaubt sind
                            % http://tex.stackexchange.com/questions/102173/


% ~~~~~~~~ Verschiedenes ~~~~~~~~~~~~~~~~~~~~~~~~~~~~~~~~~~~~~~~~~~~~~~~~~~~~~~~~~~~~~~~~~~~~~~~~~~~
\usepackage{datetime}               % typeset the month name on title page, load after babel!
% \usepackage{cite}                   % Für Mehrfachzitate wie [1-3], sonst nicht benötigt
% \usepackage{csquotes}
% \usepackage{todonotes}              % Bessere (?) Notes, als mit \marginline{...}
% \makeatletter                       % Problem with externalize:
% \renewcommand{\todo}[2][]{\tikzexternaldisable\@todo[#1]{#2}\tikzexternalenable}
% \makeatother                        % http://tex.stackexchange.com/questions/42486/
% \usepackage{blindtext}              % Für Blindtext zum testen, sonst deaktivieren!


% ~~~~~~~~ MWT-Pakete ~~~~~~~~~~~~~~~~~~~~~~~~~~~~~~~~~~~~~~~~~~~~~~~~~~~~~~~~~~~~~~~~~~~~~~~~~~~~~~
\ifcase \LANGUAGE
\else
    \usepackage{bibgerm}            % Falls 'gerunsrt' verwendet werden soll
\fi
% \usepackage{pdflscape}            % Option macht auf Institutsrechner mit externalize Probleme

\usepackage{listings}
\lstset{%
    language=Matlab,                % choose the language of the code
    basicstyle=\scriptsize,         % the size of the fonts that are used for the code
    numbers=none,                   % where to put the line-numbers
    numberstyle=\footnotesize,      % the size of the fonts that are used for the line-numbers
    stepnumber=1,                   % the step between two line-numbers. If it's 1 each line will be numbered
    numbersep=8pt,                  % how far the line-numbers are from the code
    backgroundcolor=\color{white},  % choose the background color. You must add \usepackage{color}
    showspaces=false,               % show spaces adding particular underscores
    showstringspaces=false,         % underline spaces within strings
    showtabs=false,                 % show tabs within strings adding particular underscores
    frame=false,                    % adds a frame around the code
    tabsize=3,                      % sets default tabsize to 2 spaces
    captionpos=b,                   % sets the caption-position to bottom
    breaklines=true,                % sets automatic line breaking
    breakatwhitespace=false,        % sets if automatic breaks should only happen at whitespace
    commentstyle=\color{PineGreen}, % OliveGreen, PineGreen?
    stringstyle=\color{Purple},
    escapeinside={\%*}{*)}          % if you want to add a comment within your code
}

% \usepackage{myunits}                % if siunitx is not used, e.g. 42\Ohm
% \usepackage{mymacros}               % abbrevations

% % Writing floating point numbers in german and math mode (P. Schuh)
% % without additional {}: e.g. $4,56$
% \mathchardef\ocomma="013B
% \mathchardef\pcomma="613B
% \mathcode`,="8000
% {\catcode`,=\active
% \gdef,{\obeyspaces\futurelet\next\smartcomma}}
% \def\smartcomma{\if\space\next\pcomma\else\ocomma\fi}

%% --------------------------------------------------------------------------------------------------
%% Better dealing with floats
%% Alter some LaTeX defaults for better treatment of figures:
%% See p.105 of "TeX Unbound" for suggested values.
%% See pp. 199-200 of Lamport's "LaTeX" book for details.
%%   General parameters, for ALL pages:
%\renewcommand{\topfraction}{0.9}        % max fraction of floats at top
%\renewcommand{\bottomfraction}{0.8}     % max fraction of floats at bottom
%%   Parameters for TEXT pages (not float pages):
%\setcounter{topnumber}{1}
%\setcounter{bottomnumber}{1}
%\setcounter{totalnumber}{4}     % 2 may work better
%\setcounter{dbltopnumber}{2}    % for 2-column pages
%\renewcommand{\dbltopfraction}{0.9}     % fit big float above 2-col. text
%\renewcommand{\textfraction}{0.07}      % allow minimal text w. figs
%%   Parameters for FLOAT pages (not text pages):
%\renewcommand{\floatpagefraction}{0.7}  % require fuller float pages
%% N.B.: floatpagefraction MUST be less than topfraction !!
%\renewcommand{\dblfloatpagefraction}{0.7} % require fuller float pages
%% remember using [htp] or [htpb] for placement
%% --------------------------------------------------------------------------------------------------

% ~~~~~~~~ Layout ~~~~~~~~~~~~~~~~~~~~~~~~~~~~~~~~~~~~~~~~~~~~~~~~~~~~~~~~~~~~~~~~~~~~~~~~~~~~~~~~~~
\usepackage{cmap}           % MWT
\usepackage[activate]{pdfcprot}   % character protruding (MWT)
\usepackage{setspace}       % Falls wirklich notwendig kann hiermit der Zeilenabstand geändert (MWT)
                            % werden. Nicht anders lösen, siehe LaTeX2e-Sündenregister S. 6
\setstretch{1.1}            % skip between lines (MWT)
\usepackage[                % MWT
    textwidth=160mm, textheight=230mm, top=31mm     % A4 Thesis
    % textwidth=111mm, textheight=155mm, top=25mm   % A5 Thesis
]{geometry}

% \usepackage{enumitem}       % u.a. Einzugsänderung und List-Stile ändern
% \usepackage[hyphens]{url}   % Formatieren vorn URLs per \url{http://...}
                            % SEO-URLs haben viele "-" und werden wegen Doppeldeutigkeit nicht
                            % getrennt. Um auch dort zu trennen, die Option [hyphens] übergeben
                            % Achtung: Muss zwingend vor hyperref geladen werden!
                            % http://www.froos.de/blog/latex-url-zu-lang-umbruch-erzwingen-bindestrich/
% \usepackage{float}          % Nötig, um figures per ?H? an Position festzulegen, sonst nicht.
                            % -> Wenn algorithm geladen, dann muss dieses Paket davor geladen
                            %    werden, da algorithm dieses Paket lädt. hyperref braucht ein
                            %    etwaiges float aber davor!
                            % -> tex.stackexchange.com/questions/24586/
\usepackage{hyperref}       % Links im Dokument setzen. Sollte sehr spät geladen werden!
                            % http://en.wikibooks.org/wiki/LaTeX/Hyperlinks
\hypersetup{%
    % hidelinks,            % Keine hevorgehobenen farbigen Links, erst ab Version 6.82a (05.02.2011)
}
% \hypersetup{              % Vorlage für PDF-Eigenschaften, aufnehmen in main.tex
% 	pdftitle={Praktikumsbericht},
% 	pdfsubject={Segmentierung von Laserdaten},
% 	pdfkeywords={Keyword1,Keyword2}
% }

% \usepackage[%
%   style=alphabetic,%
%   sorting=anyt,%
%   sortcites=false,%
%   sortlocale = de,%
%   firstinits=false,%
%   useprefix=false,%
%   uniquename=init,hyperref=auto,%
%   doi=false,isbn=false,url=false,%
%   minnames=3,maxnames=5,%
%   minitems=3,maxitems=99,%
%   minalphanames=3,maxalphanames=3,%
%   maxcitenames=99,%
%   autopunct=false,%
%   backend=biber%
% ]{biblatex}
% \renewcommand*{\labelalphaothers}{}
% \renewcommand\labelnamepunct{\addcomma\space} % https://tex.stackexchange.com/questions/230150/
% \bibliography{literatur}

% ~~~~~~~~ Algorithmen / Pseudocode ~~~~~~~~~~~~~~~~~~~~~~~~~~~~~~~~~~~~~~~~~~~~~~~~~~~~~~~~~~~~~~~~
% Übersicht über mögliche Pakete: http://en.wikibooks.org/wiki/LaTeX/Algorithms
% Welches davon verwenden? tex.stackexchange.com/questions/1375/
% Laut hyperref-Readme erst hyperref und dann algorithm laden
% \usepackage[Pseudocode]{algorithm} % Für eine Algorithmus-Umgebung: Für \label und \caption
% \usepackage{algorithmicx}     % Für die Umgebung des Pseudocodes, definiert Makros
% \usepackage{algpseudocode}    % Formatierung des Pseudocodes, gehört zu algorithmicx
                              % Jeder Pseudocode begint mit \begin{algorithmic}[lines]

% ~~~~~~~~ glossaries and acronyms  ~~~~~~~~~~~~~~~~~~~~~~~~~~~~~~~~~~~~~~~~~~~~~~~~~~~~~~~~~~~~~~~~
% \usepackage[
%     % toc,                        % Im Inhaltsverzeichnis erwähnen
%     % nonumberlist,               % do not show the pages where the entry is used
%                                   % especially do not use if alle entries are printed without \gls{}
%     % nopostdot,                  % no full stop at the end
%     % nogroupskip,                % do not use a gap between the letter groups
%     % xindy,                      % sort with xindy, available with TeX Live
%     % nomain,                     % Falls nicht das Hauptverzeichnis mit 'glo'-Datei verwendet wird
% ]{glossaries}                   % load AFTER hyperref (as mentioned in glossaries manual)
% \setacronymstyle{long-short}
% % \setglossarysection{section}% Festlegen, wie Glossar formatiert wird, Standard ist 'chapter'
% % \newglossary[<log-ext>]{<name>}{<in-ext>}{<out-ext>}{<title>}[<counter>]
% % Quelle: http://www.dickimaw-books.com/latex/thesis/html/makeglossaries.html#sec:newglossary
% % \newglossary[slg]{symblist}{syi}{syg}{Symbolverzeichnis}
% % \newglossary[klg]{konvlist}{kvi}{kvo}{Konventionen}
% \makeglossaries             % Erstellung der externen Dateien für die Indexierung
% % \loadglsentries{symbole}    % Vor '\begin{document}' Symbole laden, aufnehmen in main.tex


%%% +++---~~~ Ende: Pakete laden ~~~---+++ %%%%%%%%%%%%%%%%%%%%%%%%%%%%%%%%%%%%%%%%%%%%%%%%%%%%%%%%%


%%%%%%%%%%%%%%%%%%%%%%%%%%%%%%%%%%%%%%%%%%%%%%%%%%%%%%%%%%%%%%%%%%%%%%%%%%%%%%%%%%%%%%%%%%%%%%%%%%%%
%%% +++---~~~~ Ende: Einstellungen Dokumentenvorlage: header.tex ~~~---+++ %%%%%%%%%%%%%%%%%%%%%%%%%
%%%%%%%%%%%%%%%%%%%%%%%%%%%%%%%%%%%%%%%%%%%%%%%%%%%%%%%%%%%%%%%%%%%%%%%%%%%%%%%%%%%%%%%%%%%%%%%%%%%%
                                  % load the required packages and settings
\newcommand{\MYNAME}{Fabian Roos}               % your name
\newcommand{\SUPERVISOR}{Fabian Roos, M.Sc., Dipl.-Ing.\ Gandalf der Graue} % name of your supervisor
\newcommand{\FIRSTEXAMINER}{Prof.~Dr.-Ing.~Christian Waldschmidt} % your first examiner
\newcommand{\SECONDEXAMINER}{Prof.~Dr.-Ing.~Erik der Rote} % your second examiner
\newcommand{\TITLE}{%                           % this is the title for the title page
  My Final Thesis With a Very Long Title\\      % use linebreaks for optimal wrap
  Resulting Even in a Second Line\\
  Using Dynamic Space to Fill the Page
}

% uncomment and change the following lines for the PDF properties
% \hypersetup{%                                   % define the pdf settings
%   pdfauthor={Fabian Roos},
%   pdftitle={Art der Arbeit},
%   pdfsubject={Längerer Titel der Arbeit},
%   pdfkeywords={keyword1, keyword2}
% }

% \loadglsentries{inhalt/symverz}               % load the glossaries file


%%% +++---~~~ LaTeX: document beginning ~~~---+++ %%%%%%%%%%%%%%%%%%%%%%%%%%%%%%%%%%%%%%%%%%%%%%%%%%
\begin{document}
\pagenumbering{Roman}                           % first page numbers in roman number format


%%% +++---~~~ Titelseite, Sperrverker, Selbstständigkeitserklärung ~~~---+++ %%%%%%%%%%%%%%%%%%%%%%%
\thispagestyle{empty}
\def\unilogohead{%
   \hspace*{-2.2mm}%
   \parbox[b][35mm][t]{\textwidth}{%
      \parbox[t][32mm][t]{\textwidth}{%
      \includegraphics{logos/logo_brief_45_sw}%
      }%
   }%
}%

\begin{center}
\sffamily
\unilogohead
\vspace{10mm}
\Huge\textbf{\TITLE}
\vspace{10mm}

\Large\textbf{%
  \ifcase\LANGUAGE0\ifcase\TYPE0 Bachelor's Thesis\else Master's Thesis\fi
  \else\ifcase\TYPE0 Bachelorarbeit\else Masterarbeit\fi\fi
}
\vspace{20mm}

\huge\textbf{%
\MYNAME%
}
\vfill

{%
\normalsize
\monthname~\the\year%
}

\vfill

\normalsize\textbf{%
\begin{tabular}{rl}
 \ifcase\LANGUAGE0 Supervisor:\else Betreuer:\fi & \SUPERVISOR\\[1ex]
1. \ifcase\LANGUAGE0 Examiner:\else Prüfer:\fi & \FIRSTEXAMINER\\[1ex]
2. \ifcase\LANGUAGE0 Examiner:\else Prüfer:\fi & \SECONDEXAMINER%
\end{tabular}
}
\vspace{15mm}

{\small
\ifcase\LANGUAGE0
Institute of Microwave Engineering\\
Ulm University
\else
Institut für Mikrowellentechnik\\
Universität Ulm
\fi
}

\end{center}


\clearpage
                            % Titelseite laden
\cleardoubleemptypage                           % neue rechte Seite, leere Seite mit Pagestyle empty

% ...oder auch weitere Seiten einbinden

\newcommand{\DECC}{%
Ulm, {\ifcase \LANGUAGE 0 \selectlanguage{english}\else\selectlanguage{ngerman}den \fi\today}
}

%%%%%%%%%%%%%%%%%%%%%%%%%%%%%%%%%%%%%%%%%%%%%%%%%%%%
\newcommand{\DECLARATION}{%
\chapter*{\DECA}
\vspace{2\baselineskip}

\DECB
\vspace{5\baselineskip}

\DECC\hfill\MYNAME
}

\ifcase \LANGUAGE 0
\newcommand{\DECA}{Declaration}
\newcommand{\DECB}{%
I declare that I performed the presented research thesis
%
\begin{center}
\textbf{\large \TITLE}
\end{center}
%
myself and
that I wrote the thesis myself. Any additional sources of information have been
duly identified and referenced. Further, I have regarded the rules of 
scientific work of the Ulm University.
}
\else  %%%%%%%%%%%%%%%%%%%%%%%%%%%%%%%%%%%%%%%%%
\newcommand{\DECA}{Erkl\"{a}rung}
\newcommand{\DECB}{%
Hiermit erkl\"{a}re ich, die vorliegende Arbeit
%
\begin{center}
\textbf{\large \TITLE}
\end{center}
%
selbstst\"{a}ndig verfasst, keine anderen als die angegebenen Quellen
und Hilfsmittel benutzt und die Grunds\"{a}tze und Empfehlungen
\glqq Verantwortung in der Wissenschaft\grqq{} der Universit\"{a}t Ulm
beachtet zu haben.
}
\fi
\DECLARATION
\thispagestyle{empty}

\cleardoubleemptypage                           % neue rechte Seite, leere Seite mit Pagestyle empty


%%% +++---~~~ Define Abstract ~~~---+++ %%%%%%%%%%%%%%%%%%%%%%%%%%%%%%%%%%%%%%%%%%%%%%%%%%%%%%%%%%%
% If you do not know what an abstract is, you don't need it!
% -> if used, activate the abstract in the header for correct formatting
% \begin{abstract}
%   Hier steht die Zusammenfassung des Berichts, realisiert über die \emph{abstract}-Umgebung. Als
%   Überschrift steht dort \emph{Zusammenfassung}.
% \end{abstract}
% \cleardoubleemptypage                           % neue rechte Seite, leere Seite mit Pagestyle empty


%%% +++---~~~ Inhaltsverzeichnis einfügen ~~~---+++ %%%%%%%%%%%%%%%%%%%%%%%%%%%%%%%%%%%%%%%%%%%%%%%
\tableofcontents

\cleardoubleemptypage                           % Damit Seite abgeschlossen ist und die
                                                % Seitennummerierung wieder geändert werden kann


%%% +++---~~~ Dokumentbeginn: Nach Titelseite und Verzeichnisse~~~---+++ %%%%%%%%%%%%%%%%%%%%%%%%%%%
\pagenumbering{arabic}		                    % Ab dem eigentlichen Inhalt
                                                % arabische Numerierung von 1 an

% Inhalte einbinden
\chapter{Basic Mathematics}
My first \LaTeX\ example with footnote for abc\footnote{means bac}. And this are my first formulae:
\begin{align}
  \SI{10}{\ohm} - \SI{9}{\ohm} &= \SI{1}{\ohm}\\
  % 10\Ohm -9\Ohm &= 1\Ohm % with myunits which must be activated in the header
  a^2 + b^2 &= c^2\,,\nonumber\\
  \leadsto c &= \sqrt{a^2+b^2}\label{eq:pythagoras}\,,\\
  a &= \operatorname{e}^{-\mathrm{i}\omega t}\,,\\
  a &= \frac{1}{2}\,,\\
  a_{n+1} &= a_n^2+\sum_{i=0}^{n}b^{2i}\,,\\
  \intertext{with text inbetween without changing the alignment of the equations}
  a_{n'+1}' &= a_{n'}^2+\sum_{i=0}^{n'}b^{2i}\,,\\
  \boldsymbol{b} &= \boldsymbol{a }+ \boldsymbol{c}\,,\\
  a &= a\cdot\, i\,,\qquad i\in\mathbb{N}\,.\\
  \int_a^b (x+x^2)\mathrm{d} x &= \int_a^b (x+x^2)\mathrm{d} x\\
  \oint_a^b (x+x^2)\mathrm{d}x &= \int_a^b (x+x^2)\mathrm{d}x\\
  \oiint_a^b (x+x^2)\mathrm{d}x &= \int_a^b (x+x^2)\mathrm{d}x\\
  \iiint_a^b (x+x^2)\mathrm{d}x &= \int_a^b (x+x^2)\mathrm{d}x\\
  \int\int\int_a^b (x+x^2)\mathrm{d}x &= \int_a^b (x+x^2)\mathrm{d}x\\
  \stackrel{\leftrightarrow}{A} &= \stackrel{\leftrightarrow}{B}+\stackrel{\leftrightarrow}{C}\\
  R_{\text{max}}& = \frac{T_{\mathrm{c}} \cdot c}{2 B} f_{\mathrm{max}}
\end{align}

$\pi$ is roughly\ldots

It's obvious, that Pythogoras' equation~(\ref{eq:pythagoras}) on page~\pageref{eq:pythagoras} is
wrong. The correct relation between the three lines in a rectangular triangle (without equation
number) is:
\begin{align*}
  c &= \sqrt[2.1]{a^2+b^2}\,.
\end{align*}

\begin{table}[!b]
  \centering
  \begin{tabular}{lr}
    Parameter                               & Value\\
    \hline
    carrier frequency $f_{\text{c}}$        & \SI{77}{\giga\hertz}\\
    bandwidth $B$                           & \SI{500}{\mega\hertz}\\
    chirp duration $T_{\text{c}}$           & \SI{20}{\micro\second}\\
    chirp repetition time $T_{\text{r}}$    & \SI{25}{\micro\second}\\
    sampling frequency $f_{\text{s}}$       & \SI{50}{\mega\hertz}\\
    number of chirps $L$                    & 128\\
    window function                         & Hann\\
    zero padding                            & signal length doubled
  \end{tabular}
  \caption{Specifications of Radar Parameters.}
  \label{tab:radar_para}
\end{table}

\clearpage % use a clearpage carefully and only if you know the sideeffects
\label{ici}
We are now on page number~\pageref{ici}.

\chapter{TikZ and PGFPlots example}
Example for the \emph{subcaption} package using a TikZ and PGFPlots image.

\begin{figure}[tbp]
  \centering
  \subcaptionbox{This is a TikZ image which can be scaled only with a fixed axis ratio.
  \label{fig:tikzimg}}
  [0.47\linewidth]{\includegraphics[width=0.47\linewidth]{gfx/chirp_sequence_shift_01.tikz}}
  \hfill
  \subcaptionbox{This is a PGFPlots example. Take a look at the source code for further examples.
    Using the tikzscale package the axis ratio can be adepted.
  \label{fig:pgfplotsimg}}
  [0.47\linewidth]{\includegraphics[width=0.47\linewidth,height=0.2\textheight]{gfx/example_plots.tikz}}
  \caption{This is an example for the subcaption package.}
  \label{fig:subcapex}
\end{figure}

If in image~\ref{fig:subcapex} a linewidth smaller than 0.49 is chosen, the images are positioned
left and right. If the size is larger as now, they are positioned differently. A reference to
Bild~\ref{fig:subcapex}~(\subref{fig:tikzimg}). Also note that the axis ratio for a PGFPlots image
can be influenced.

\begin{figure}[tbp]
  \centering
  \subcaptionbox{This is a TikZ image which can be scaled only with a fixed axis ratio.
  \label{fig:tikzimg2}}
  [0.60\linewidth]{\includegraphics[width=0.6\linewidth]{gfx/chirp_sequence_shift_01.tikz}}
  \par\bigskip
  \subcaptionbox{This is a PGFPlots example. Take a look at the source code for further examples.
    Using the tikzscale package the axis ratio can be adepted.
  \label{fig:pgfplotsimg2}}
  [0.60\linewidth]{\includegraphics[width=0.6\linewidth,height=0.35\textheight]{gfx/example_plots.tikz}}
  \caption{This is an example for the subcaption package.}
  \label{fig:subcapex2}
\end{figure}

The examples in Bild~\ref{fig:subcapex} and~\ref{fig:subcapex2} use either a TikZ image or a
PGFPlots image consisting of a function or data points. For large data matrices this approach is not
useful. Here use the \emph{graphics} command of PGFPlots. See~\ref{fig:pgfplots} for an example.

\begin{figure}[tbp]
  \begin{minipage}[b][\textheight][b]{\linewidth}
    \centering
    \subcaptionbox{This is an image which was created using the \emph{imagesc} plot function from
      Matlab in combination with matlab2tikz.
    \label{fig:imagesc}}
    [\linewidth]{\includegraphics[width=0.95\linewidth]{gfx/surfimgsc.tikz}}
    \vfill
    \subcaptionbox{For 3D plots there is no suitable way with matlab2tikz. Create the png on your own
      and create an axis environment with PGFPlots.
    \label{fig:surf2tikz3d}}
    [\linewidth]{\includegraphics[width=0.65\linewidth]{gfx/surf2tikz3D.tikz}}
    %
    \caption{An example for 2D and 3D plots.}
    \label{fig:pgfplots}
  \end{minipage}
\end{figure}

\begin{figure}[tbp]
  \centering
  \tikzset{%
	% Self defined bulding blocks. 
	% Nevertheless circutikz has implemented filters, couplers and other components since version 0.4, they are mostly implemented as bipoles.
	% The usage of bipoles: \draw (start) to[lowpass/amp/adc,....] (end).
	% The problem is, that if one wants to use arrows, the arrows in bipoles can not be sat manual (fixed in circuitikz source) AND THEY ARE NOT CONSISTENT
	% Also it's quite a mess, which component is a monopole, simple block, bipol, quad/triple etc
	% Following are a few examples on how to define your own blocks. 
	%
	% % % % % % % % % % % % % % % % % % % % % % % % % % % % % % % % % % % % % % % % % % % % % % % % % % % % % % % % % % % % % % % % % % % % %
	% % % % % % % % % % % % % % % % % % % % % % % % % % % % % % % % % % % % % % % % % % % % % % % % % % % % % % % % % % % % % % % % % % % % %
	% % % % % % % % % % % % % % % % % % % % % % % % % % % % % % % % % % % % % % % % % % % % % % % % % % % % % % % % % % % % % % % % % % % % %
	% % % % % % % % % % % % % % % % % % % % % % % % % % % % % % % % % % % % % % % % % % % % % % % % % % % % % % % % % % % % % % % % % % % % %
	%
	% Standard block definition, the width and height is adopted from the circutizk source code, so don't mind the strange values. Also the linewidth is set according to the circutrikz source code.
	block/.style    	= 	{draw, fill=white, thick, rectangle, minimum height = 0.98cm, minimum width = 0.98cm, node distance=2.5cm, line width=1.5pt},
	%
	% Standard circular block
	circleblock/.style	= 	{draw, fill=white, thick, circle, minimum width = 0.98cm,  line width=1.5pt, node distance=2.5cm},
	%
	% Label for circuitikz nodes, as they're reference is in the middle and not on the outer edge of the node....
	circuitikzlabel/.style	=	{label={[label, label distance=0.5cm]#1}},
	%
	%
	%
	% VCO/Oscillator 
	myVCO/.style			= 	{circleblock, path picture={%
		\draw[line width=0.75pt] 	($(path picture bounding box.west)+(0.09cm,0)$) sin ($(path picture bounding box.center)-(0.2cm,-0.2cm)$) cos  (path picture bounding box.center) sin ($(path picture bounding box.center)-(-0.2cm,0.2cm)$) cos ($(path picture bounding box.east)-(0.09cm,0)$);
		}
	},
	% Amplifier, as circuitikz does only provite amplifiers as 2-ports/bipoles
	myAMP/.style		= 	{block, node distance=2.5cm, path picture={%
		\draw[fill=white, line width=0.75pt] ($(path picture bounding box.center)+(0.7em,0)$) -- ($(path picture bounding box.center)-(0.7em,-0.7em)$) -- ($(path picture bounding box.center)-(0.7em,0.7em)$)  -- cycle;
		}
	},
	% Same for ADC
	myADC/.style 	=	{block, path picture={%
		\draw[line width=0.75pt] 	(path picture bounding box.south west) -- (path picture bounding box.north east);
		\node[] at ($(path picture bounding box.center)+(-.5em,.5em)$) () {D};
		\node[] at ($(path picture bounding box.center)+(.5em,-.5em)$) () {A};
		} 
	},
	% Same for filters
	myLP/.style	=	{block, path picture={%
		%Sine-Waves
		\draw[line width=.75pt] 	($(path picture bounding box.west)+(0.3em,0)$) sin ($(path picture bounding box.center)-(0.50em,-0.3em)$) cos  (path picture bounding box.center) sin ($(path picture bounding box.center)-(-0.50em,0.3em)$) cos ($(path picture bounding box.east)-(0.3em,0)$);
		\draw[line width=0.75pt] 	($(path picture bounding box.west)+(0.3em,-0.65em)$) sin ($(path picture bounding box.center)-(0.50em,0.35em)$) cos  ( $(path picture bounding box.center)-(0,0.65em)$) sin ($(path picture bounding box.center)-(-0.50em,0.95em)$) cos ($(path picture bounding box.east)-(0.3em,0.65em)$);
		\draw[line width=0.75pt] 	($(path picture bounding box.west)+(0.3em,0.65em)$) sin ($(path picture bounding box.center)-(0.50em,-0.95em)$) cos  ( $(path picture bounding box.center)+(0,0.65em)$) sin ($(path picture bounding box.center)-(-0.50em,-0.35em)$) cos ($(path picture bounding box.east)-(0.3em,-0.65em)$);
		% Cancelation
		\draw[line width=0.75pt] 	($(path picture bounding box.center)-(0.2em,0.2em)$) -- (path picture bounding box.center) -- ($(path picture bounding box.center)+(0.2em,0.2em)$) ;
		\draw[line width=0.75pt] 	($(path picture bounding box.center)-(0.2em,-0.45em)$) -- ($(path picture bounding box.center)+(0,0.65em)$) -- ($(path picture bounding box.center)+(0.2em,0.85em)$) ;
		}
	},
}
\begin{tikzpicture}[line width=0.7pt,>=latex,node distance=2.5cm]
	% First: All building blocks are placed relative to the first component
	\draw (0,0)
		node[myVCO, label={below:VCO}] (oszi) {}
		% As the coupler ports are not in the middle, based on the size (again extraceted from circutikz source code), an yshift is perfomed to have the input on the same height as the output of the oszillator. The xshift is used to place the VCO and ADC on the same y-value after all.
		node[coupler, right of=oszi, yshift=-0.455cm, xshift=1cm] (coupler) {}
		% Undo the yshift as the output of the coupler is on the same height as the input of the amplifier
		node[myAMP, right of=coupler, yshift=0.455cm, node distance=4cm, label={below:PA}] (pa) {}
		% Circulator is rotated that the ports are on the correct position, normally ports are arranged as follows:
		% 
		%				  o------------o
		%						|
		%						o
		node[circulator, below right of=pa, rotate=90] (circ) {}
		node[antenna, right of = circ] (antenna) {}
		node[myAMP, below left of=circ, rotate=180, label={LNA}] (lna) {}
		% Used redefinition of label, otherwise the label would be overlapping with the mixer shape
		node[mixer, left of=lna, circuitikzlabel={below:mixer}] (mixer) {}
		node[myLP, left of=mixer, label={below:lowpass}] (lowpass) {}
		node[myADC, left of=lowpass] (adc) {};
	
	% Connect everything together
	\draw[->] (oszi) -- (coupler.4);
	\draw[->] (coupler.3) -- (pa.west);
	% Match is placed relative to coupler port, yscale=-1 mirrors the component at the y-axis
	\draw[] (coupler.1) -| ++(-0.5cm,-0.1cm) node[match, rotate=-90, yscale=-1] {};
	\draw[->] (pa.east) -| (circ.2);
	\draw[-] (circ.3) -- (antenna);
	\draw[->] (circ.1) |- (lna.west);
	\draw[->] (lna.east) -- (mixer.east);
	\draw[->] (coupler.2) -| (mixer.north);
	\draw[->] (mixer.west)  -- (lowpass.east);
	\draw[->] (lowpass.west) -- (adc.east);
\end{tikzpicture}

  \caption{A \emph{circuitikz} example.}
  \label{fig:cwradar}
\end{figure}

\begin{figure}[tbp]
  \centering
  \includegraphics[width=\linewidth]{gfx/smith_chart.tikz}
  \caption{A Smith Chart example.}
  \label{fig:sc}
\end{figure}

\begin{figure}[tbp]
  \centering
  \includegraphics{gfx/fc_example.tikz}
  \caption{A flow chart example.}
  \label{fig:fc}
\end{figure}

If you use several TikZ oder PGFPlots consider to active the \emph{externalize} option. Then think about a
shell-escape.

\chapter{Include literature}
As an example a literature is included. Only entries which are referenced as like~\cite{Roos2015}
are printed.



% Testinhalt laden, dafür muss Paket 'blindtext' in 'header.tex' geladen werden
%\blinddocument % kleines Dokument
%\Blinddocument % großes Dokument


%% Anhang einleiten
%\appendix


%% Verzeichnisse
%\chapter*{Symbole}
%\addcontentsline{toc}{chapter}{Symbole}
%\glsaddall[types={symblist}]
%\glsaddall[types={konvlist}]
%\printglossary[type=konvlist, style=long]
%\printglossary[type=symblist,style=long]

%\listoffigures
%\listoftables

\bibliography{literatur}
% \printbibliography

\ifcase \LANGUAGE
    % \bibliographystyle{unsrt}                   % do not use 'unsrt' anymore, due to compatibility issues
    \bibliographystyle{bibtex/IEEEtran}         % use the IEEE default
    % \bibliographystyle{bibtex/IEEEtranSA}       % a IEEE style but with alphanumeric behaviour
\else
    \bibliographystyle{bibtex/myIEEEtran-de}    % use the IEEE default, german adaptions by Tobias Chaloun
    % \bibliographystyle{alphadin}
    % \bibliographystyle{unsrtdin}
    % \bibliographystyle{gerunsrt}                % do not use 'gerunsrt' anymore, due to compatibility issues
\fi

\end{document}
